\documentclass[11pt]{article}
\usepackage[margin=0.9in]{geometry}
\usepackage{amsmath,amssymb,amsthm}
\usepackage[hidelinks]{hyperref}
\newtheorem{theorem}{Theorem}
\newtheorem{proposition}[theorem]{Proposition}
\newtheorem{lemma}[theorem]{Lemma}
\newtheorem{corollary}[theorem]{Corollary}
\newcommand{\F}{\mathbb F}
\newcommand{\ad}{\operatorname{ad}}
\newcommand{\tr}{\operatorname{tr}}
\newcommand{\GL}{\operatorname{GL}}
\newcommand{\slie}{\mathfrak{sl}}
\title{Explicit Lie polynomials with prescribed images over finite fields}
\author{Henry Zweiman}
\date{September 9, 2026}
\begin{document}
\maketitle
\begin{abstract}
For every odd prime power $q$, we give a closed formula in two Lie variables
whose image on $\slie_2(\F_q)$ is any prescribed conjugation-invariant subset
containing zero. The degree is at most $6q-5$. The construction converts an
arbitrary scalar function of the quadratic invariant into a Lie polynomial
on the locus where the commutator has nonzero determinant, and vanishes
elsewhere. It answers the explicit-construction question of Kishnani and
Singh in rank one. We determine all fibers of the resulting class selectors:
each selected nonzero conjugacy class receives the same number of input
pairs. Finally, any Lie polynomial whose image is zero together with one
semisimple class has degree at least $\lceil q/4\rceil$, even when arbitrarily
many variables are allowed. Thus a linear degree bound has the correct order.
\end{abstract}

\section{Introduction and the explicit formula}
Let $F=\F_q$, where $q$ is an odd prime power, and put $L=\slie_2(F)$.
We use the bracket $[X,Y]=XY-YX$ and the convention
$\ad_X(Y)=[X,Y]$. A Lie polynomial belongs to the free Lie algebra over $F$;
its syntax permits scalar coefficients, addition, and brackets, but no
nonzero matrix constants. Its total degree is the largest number of variable
occurrences in a nonzero homogeneous component.

Lie-polynomial images are invariant under conjugation and contain zero.
Lubotzky's theorem \cite{Lubotzky} gives an analogous image classification
for word maps on finite simple groups. Kishnani and Singh \cite[Theorem 1.3]{KS}
prove an existence theorem for finite simple Chevalley Lie algebras.
Their Question 1.2 asks for explicit realizing polynomials; their Section 9
gives partial constructions for $\slie_2(\F_q)$.
The contribution here is a uniform two-variable formula for every admissible
subset in this rank, together with a degree bound and its exact distribution.
The existence classification itself is prior work.

For $X\in L$ define
\[
 Q(X)=-\det X,\qquad
 B(X,Y)=\tfrac12\tr(XY),\qquad
 \Delta(X,Y)=Q([X,Y]).
\]
These invariants are used to describe evaluations. They do not appear as
operations in the Lie polynomial below. For
$r(t)=\sum_{k=0}^{q-1}r_kt^k\in F[t]$, set
\begin{equation}\label{eq:word}
 W_r(X,Y)=-\sum_{k=0}^{q-1}r_k4^{-k}
       \ad_{[X,Y]}^{\,2q-3}\bigl(\ad_X^{\,2k+2}(Y)\bigr).
\end{equation}
Here $4^{-k}$ denotes an element of $F$; in particular the formula is valid
in characteristic three. Each summand is an iterated bracket in $X,Y$.

\newpage
\begin{theorem}\label{thm:main}
For every polynomial $r\in F[t]$ of degree at most $q-1$,
\begin{equation}\label{eq:evaluation}
 W_r(X,Y)=\Delta(X,Y)^{q-1}r(Q(X))X
 \qquad(X,Y\in L).
\end{equation}
For $T\subseteq F$ define
\begin{equation}\label{eq:indicator}
 r_T(t)=\sum_{a\in T}\bigl(1-(t-a)^{q-1}\bigr),\qquad W_T=W_{r_T}.
\end{equation}
Then
\begin{equation}\label{eq:image}
 W_T(L^2)=A_T:=\{0\}\cup\{X\in L\setminus\{0\}:Q(X)\in T\}.
\end{equation}
Every conjugation-invariant subset of $L$ containing zero is $A_T$ for a
unique $T\subseteq F$. Consequently every such subset is the image of the
explicit Lie polynomial \eqref{eq:word}, of degree at most $6q-5$.
\end{theorem}

The theorem gives a constructive answer to \cite[Question 1.2]{KS} for
$L=\slie_2(F)$, with no restriction on the odd prime power $q$ or on the
chosen union. In particular it permits arbitrary mixtures of split,
nonsplit, and nilpotent classes. It does not assert an explicit construction
in higher rank.

\section{The invariant identity}
We prove the matrix identities over any field of characteristic different
from two, before using finiteness.

\begin{lemma}\label{lem:identities}
For $X,Y\in\slie_2(F)$ and $C=[X,Y]$, one has
\begin{align}
 X^2&=Q(X)I, & XY+YX&=2B(X,Y)I,\label{eq:CH}\\
 \ad_X^2(Y)&=4\bigl(Q(X)Y-B(X,Y)X\bigr),\label{eq:double}\\
 B(X,C)&=0, & \ad_C^2(X)&=4Q(C)X,\label{eq:orth}\\
 [C,\ad_X^2(Y)]&=-4Q(C)X.\label{eq:transfer}
\end{align}
Moreover, for every $k\geq0$,
\begin{equation}\label{eq:recurrence}
 \ad_X^{2k+2}(Y)=(4Q(X))^k\ad_X^2(Y).
\end{equation}
\end{lemma}

\begin{proof}
The first equality of \eqref{eq:CH} is Cayley--Hamilton for a trace-zero
matrix. Polarizing it gives the second, since
$Q(X+Y)-Q(X)-Q(Y)=\tr(XY)$.
Multiplication by $X$ gives $XYX=2B(X,Y)X-Q(X)Y$, whence
\[
 [X,[X,Y]]=X^2Y-2XYX+YX^2
          =4Q(X)Y-4B(X,Y)X.
\]
Cyclicity of trace gives
$2B(X,C)=\tr(X^2Y)-\tr(XYX)=0$.
Applying \eqref{eq:double} to $C,X$ proves \eqref{eq:orth}.
As $\ad_X^2(Y)=[X,C]$, antisymmetry now gives
\[
 [C,\ad_X^2(Y)]=[C,[X,C]]=-[C,[C,X]]=-4Q(C)X.
\]
Finally, applying $\ad_X$ to \eqref{eq:double} gives
$\ad_X^3(Y)=4Q(X)\ad_X(Y)$, since $\ad_X(X)=0$.
Iteration proves \eqref{eq:recurrence}; no division by $Q(X)$ is used.
\end{proof}

\begin{proposition}\label{prop:lift}
For each $k\geq0$ and odd prime power $q$,
\begin{equation}\label{eq:monomial}
 -4^{-k}\ad_C^{2q-3}\bigl(\ad_X^{2k+2}(Y)\bigr)
       =Q(C)^{q-1}Q(X)^kX
       \quad(C=[X,Y]).
\end{equation}
\end{proposition}

\begin{proof}
Write $\Delta=Q(C)$. The exponent $2q-3$ is $1+2(q-2)$.
By \eqref{eq:recurrence}, \eqref{eq:transfer}, and \eqref{eq:orth}, the
left side of \eqref{eq:monomial} is
\begin{align*}
 -4^{-k}(4Q(X))^k\ad_C^{2(q-2)}(-4\Delta X)
 &=4^{q-1}\Delta^{q-1}Q(X)^kX\\
 &=\Delta^{q-1}Q(X)^kX.
\end{align*}
The final equality uses $4^{q-1}=1$ in $F$.
All preceding identities hold when either invariant is zero. Thus no
exceptional locus was removed during the computation.
\end{proof}

Summing \eqref{eq:monomial} proves \eqref{eq:evaluation}. This also gives a
useful formulation independent of image selection: for any scalar function
$f:F\to F$, take its unique interpolation polynomial $r$ of degree less
than $q$. Formula \eqref{eq:word} realizes the map $X\mapsto f(Q(X))X$
whenever $\Delta(X,Y)\neq0$, and returns zero otherwise.
Only the coefficients of the scalar interpolation polynomial need to be
computed; no search through free-Lie words is required.

\section{Conjugacy classes and proof of surjectivity}
\begin{lemma}\label{lem:companion}
For each $t\in F$, the set
$\{X\in L\setminus\{0\}:Q(X)=t\}$ is one $\GL_2(F)$-conjugacy class.
Every nonzero $X\in L$ has a partner $Y\in L$ satisfying
$\Delta(X,Y)=1$.
\end{lemma}

\begin{proof}
A nonzero trace-zero matrix is not scalar, since the characteristic is odd.
It therefore has a cyclic vector $v$: otherwise every vector would be an
eigenvector, and applying this to two independent vectors and their sum
would make the matrix scalar. In the basis $(Xv,v)$, its matrix is
\begin{equation}\label{eq:companion}
 X_t=\begin{pmatrix}0&1\\t&0\end{pmatrix},\qquad t=Q(X),
\end{equation}
since $X^2v=tv$. Thus every nonzero matrix of invariant $t$ is conjugate
to $X_t$. Conversely the invariant is preserved by conjugation.
With
\[
 Y_0=\begin{pmatrix}0&0\\1&0\end{pmatrix}
 \quad\hbox{one has}\quad
 [X_t,Y_0]=\begin{pmatrix}1&0\\0&-1\end{pmatrix},
\]
whose $Q$-value is one. Conjugate $Y_0$ back with the same basis change.
\end{proof}

\begin{proof}[Proof of Theorem \ref{thm:main}]
The evaluation identity follows from Proposition \ref{prop:lift}.
For every $t\in F$, the summand $1-(t-a)^{q-1}$ is one if $t=a$ and zero
otherwise. Hence $r_T$ is the indicator of $T$, even in nonprime fields.
The same finite-field identity shows that $\Delta^{q-1}$ is the indicator
of $\Delta\neq0$. Consequently
\begin{equation}\label{eq:selection}
 W_T(X,Y)=
 \begin{cases}
 X,&Q(X)\in T\ \hbox{and}\ \Delta(X,Y)\neq0,\\
 0,&\hbox{otherwise}.
 \end{cases}
\end{equation}
Every nonzero $X\in A_T$ is attained by Lemma \ref{lem:companion}, and
zero is attained at $(0,0)$. This proves \eqref{eq:image}.
The same lemma classifies the nonzero conjugacy classes, proving the
assertion about all admissible subsets and uniqueness of $T$.
The $k$th summand of \eqref{eq:word} has degree at most
\[
 2(2q-3)+(2k+2)+1=4q+2k-3\leq6q-5.
\]
This proves the degree bound, including $T=\varnothing$, for which the
word is zero.
\end{proof}

It follows that for a subset $A\subseteq L$ the following are equivalent:
$A$ is a Lie-polynomial image; $A$ contains zero and is invariant under
$\GL_2(F)$-conjugation; and $A=A_T$ for some $T\subseteq F$.
Necessity of conjugation invariance follows by applying a conjugation to
every input of a Lie polynomial. No theorem identifying all Lie algebra
automorphisms is needed for this argument.

\section{Exact fibers and distribution on classes}
Let $\chi:F^*\to\{1,-1\}$ be the quadratic character. The classes with
$t\neq0$ are semisimple; they are split for $\chi(t)=1$ and nonsplit for
$\chi(t)=-1$. The class with $t=0$ consists of the nonzero nilpotents.

\begin{theorem}\label{thm:fibers}
Fix $T\subseteq F$ and a nonzero $Z\in A_T$, and write $t=Q(Z)$. Then
\begin{equation}\label{eq:fibers}
 |W_T^{-1}(Z)|=
 \begin{cases}
 q^2(q-1),&t=0,\\
 q(q-1)^2,&t\neq0,\ \chi(t)=1,\\
 q(q^2-1),&t\neq0,\ \chi(t)=-1.
 \end{cases}
\end{equation}
Every selected nonzero conjugacy class has exactly
$q^2(q^2-1)(q-1)$ preimages. In particular,
\begin{equation}\label{eq:zero}
 |W_T^{-1}(0)|=q^6-|T|q^2(q^2-1)(q-1).
\end{equation}
\end{theorem}

\begin{proof}
By \eqref{eq:selection}, a nonzero output $Z$ forces the first input to be
$Z$. Its fiber therefore consists precisely of the $Y$ for which
$Q([Z,Y])\neq0$. Conjugate $Z$ to $X_t$ in \eqref{eq:companion}, and write
\[
 Y=\begin{pmatrix}d&a\\b&-d\end{pmatrix}.
\]
A direct multiplication gives
\begin{equation}\label{eq:binary}
 Q([X_t,Y])=(b-ta)^2-4td^2.
\end{equation}
For each $a$, the variable $v=b-ta$ ranges independently over $F$.
If $t=0$, the zero set of \eqref{eq:binary} has $q^2$ elements.
If $t=s^2\neq0$, the equation is $(v-2sd)(v+2sd)=0$; the two distinct
lines have $2q-1$ points altogether. With the free choice of $a$, this
gives $q(2q-1)$ zeros. If $t$ is a nonsquare, $v^2=4td^2$ forces $v=d=0$,
giving $q$ zeros. Subtracting from $q^3$ proves \eqref{eq:fibers}.

For completeness, the sizes of the three types of nonzero classes are
\begin{equation}\label{eq:classsizes}
 q^2-1,\qquad q(q+1),\qquad q(q-1),
\end{equation}
respectively. These can be counted without a centralizer formula. Matrices
$\left(\begin{smallmatrix}d&a\\b&-d\end{smallmatrix}\right)$ with
$Q=t$ satisfy $ab=t-d^2$. For a fixed $d$ there are $q-1$ choices of $(a,b)$
if $t-d^2\neq0$, and $2q-1$ if it is zero. The number of roots of $d^2=t$
is one, two, or zero for the three cases, respectively. For $t=0$ remove
the zero matrix. This yields \eqref{eq:classsizes}.
Multiplying each class size by the corresponding value in \eqref{eq:fibers}
gives $q^2(q^2-1)(q-1)$ in all three cases. Finally $|L^2|=q^6$, and the
$|T|$ selected nonzero classes are disjoint, proving \eqref{eq:zero}.
\end{proof}

Thus, conditional on a nonzero output from uniformly distributed inputs,
the selected conjugacy classes are equally likely. Within any one class,
the output is uniform as well. Uniformity over the entire selected set
need not hold, since the class sizes differ.

\section{Degree and arity limitations}
The construction has degree proportional to the field size. The next
bound shows that this dependence cannot be replaced by a sublinear bound
valid for every admissible image.

\begin{proposition}\label{prop:lower}
Let $w$ be a Lie polynomial over $F$, in any finite number of variables,
whose image on $L$ is $\{0\}$ together with one nonzero semisimple conjugacy
class. If $D$ is its total degree, then $D\geq\lceil q/4\rceil$.
\end{proposition}

\begin{proof}
Let the selected class have invariant $a\neq0$. Choose an input tuple
$(X_1,\ldots,X_m)$ on which $w$ is nonzero, and introduce a scalar
indeterminate $s$. The ordinary polynomial
\[
 f(s)=Q\bigl(w(sX_1,\ldots,sX_m)\bigr)\in F[s]
\]
has degree at most $2D$, with $f(0)=0$ and $f(1)=a$.
For every $s\in F$, its value is either zero or $a$. Hence
$f(s)(f(s)-a)$ has all $q$ field elements as roots. It is a nonzero
polynomial: neither factor is identically zero, by the two displayed
values, and $F[s]$ is an integral domain. Its degree is at most $4D$,
so the root bound gives $q\leq4D$.
\end{proof}

Theorem \ref{thm:main} and Proposition \ref{prop:lower} establish the
optimal order of growth of the worst-case degree. They do not determine
the best constant or the minimum degree for an individual image.
Two variables are also necessary for every proper image other than $\{0\}$:
a free Lie algebra on one generator is one-dimensional, so its polynomial
maps are $X\mapsto cX$, with images only $\{0\}$ and $L$.

As a basic example, over $\F_3$ the nilpotent selector uses
$r_{\{0\}}(t)=1-t^2$. Since $4=1$ in this field, the formula becomes
\[
 W_{\{0\}}(X,Y)
 =-\ad_{[X,Y]}^3\bigl(\ad_X^2(Y)\bigr)
  +\ad_{[X,Y]}^3\bigl(\ad_X^6(Y)\bigr).
\]
Each of the eight nonzero nilpotents has eighteen preimages, and the zero
fiber has $3^6-8\cdot18=585$ elements. The same formula \eqref{eq:word},
with the appropriate interpolation coefficients, handles every subset of
classes over every odd extension field.

\section{Verification and further questions}
The proofs above use only matrix identities, finite-field interpolation,
and elementary counting. The accompanying exact checker independently
expands the three basic identities in generic $2\times2$ matrices over
an integer polynomial ring. It also evaluates the actual nested brackets
in \eqref{eq:word}, rather than using \eqref{eq:evaluation} to compute them.
It checks every input pair over $\F_3$ and $\F_5$, all companion first inputs
and all second inputs over $\F_7$ and $\F_9$, and a reproducible sample over
$\F_{25}$. The nonprime fields are implemented as quadratic extensions,
so their elements are not incorrectly treated as residues modulo $q$.
All forty admissible subsets over $\F_3$ and $\F_5$ have their complete
images and zero fibers checked directly. These finite tests support the
universal proof; they are not a replacement for it.

The formula suggests two further problems. First, determine the smallest
possible degree for a prescribed semisimple class together with zero.
Second, construct analogous explicit invariant selectors in higher-rank
simple Lie algebras. The rank-one proof works because a double adjoint
operation transfers a scalar invariant onto the first input, and the
quadratic invariant classifies all nonzero conjugacy classes. Neither
property is available in this form in higher rank.

\begin{thebibliography}{9}
\bibitem{KS}
H. Kishnani and A. Singh,
\emph{Images of Lie Polynomials on simple Lie algebras},
arXiv:2605.19512v1 (2026).
\url{https://arxiv.org/abs/2605.19512}.
\bibitem{Lubotzky}
A. Lubotzky,
\emph{Images of word maps in finite simple groups},
Glasgow Math. J. \textbf{56} (2014), no. 2, 465--469.
\url{https://doi.org/10.1017/S0017089513000396}.
\end{thebibliography}
\end{document}
